\section{An elementary replacement for both deep exponent pairs}
\label{app:elementarypair}

We use the spacing, thin-energy and stationary-support bounds proved
in \cite[Sections 2--3]{Thin}. Here we derive the derivative estimates
by direct differencing and the upper-amplitude bound from the curved sum.

\begin{theorem}[eighth moment using the sparse fourth moment]
\label{thm:elementaryeighth}
Using the sparse-fourth-moment theorem \cite[(\ref*{six-eq:basefourth})]{Weighted}, we obtain
\[
 \int_0^T|\zeta(1/2+it)|^8dt
 \ll_\varepsilon T^{357082/242699+\varepsilon}.
\]
\end{theorem}
The exponent-pair estimates for this statement follow from the
finite differencing arguments below. The short-second-moment estimate
is proved in \cite[Appendix~\ref*{six-app:atkinsondirect}]{Weighted}; the sparse fourth moment
is the cited analytic input.

\begin{lemma}[finite derivative tests from differencing]
\label{lem:elementaryderivatives}
Let $2\le d\le5$, let $I$ be an interval of length at most $M\ge2$,
and suppose $f\in C^d(I)$ with
$\lambda\le|f^{(d)}|\le A\lambda$. Then
\[
 \left|\sum_{m\in I}e(f(m))\right|
 \ll_{d,A} M\lambda^{1/(2^d-2)}
   +M^{1-2^{2-d}}\lambda^{-1/(2^d-2)}.
\]
Hence, if $|f^{(d)}|\asymp FM^{-d}$ for $d=3,4,5$ and
$F\ge M$, on any subinterval of a fixed dilation of $(M,2M]$, then
\begin{equation}\label{eq:elementarypair}
 \left|\sum_{m\in I}e(f(m))\right|
 \ll M^{5/6}F^{1/30}.
\end{equation}
\end{lemma}
Thus $(k,\ell)=(1/30,13/15)$, with $a=1-\ell+k=1/6$, gives
exactly the finite-derivative conditions needed by our reciprocal phases.
\begin{proof}
The case $d=2$ is \cite[Lemma~\ref*{six-lem:secondproof}]{Weighted}; its additional constant
is absorbed since the displayed two terms have sum at least a constant.
Extend the summand by zero outside $I$ and average
its first $H$ translates. Cauchy's inequality on their union, of length
$O(M)$ for $H\le M$, and expansion of the square give
\[
 |S|^2\ll M^2/H+\frac MH\sum_{1\le h<H}
 \left|\sum_{m\in I\cap(I-h)}e(f(m+h)-f(m))\right|.
\]
The factor $1-h/H$ in the exact expansion has been bounded by one.
The differenced $(d-1)$st derivative is the integral of $f^{(d)}$
over a segment of length $h$, hence has fixed sign and size
$\asymp h\lambda$ on its interval. Write
$\alpha=(2^{d-1}-2)^{-1}$ and $\beta=1-2^{3-d}$.
Induction and the elementary sums of $h^{\pm\alpha}$ yield
\[
 |S|^2\ll M^2/H+M^2\lambda^\alpha H^\alpha
                   +M^{1+\beta}\lambda^{-\alpha}H^{-\alpha}.
\]
For $0<\lambda<c$ take an integer
$H\asymp\lambda^{-\alpha/(1+\alpha)}$, when this is at most $M$.
Taking square roots gives the statement, since
$\alpha/(2+2\alpha)=(2^d-2)^{-1}$ and
$(1+\beta)/2=1-2^{2-d}$.
If the proposed $H$ exceeds $M$, then
$\lambda\ll M^{-(2^{d-1}-1)}$, and the second term in the claimed
bound is at least a constant times $M$: indeed
$1-2^{2-d}+(2^{d-1}-1)/(2^d-2)\ge1$ for $d\ge3$.
For $\lambda\ge c$, the first term gives the trivial bound.
Floors and intervals of length less than two are absorbed by constants.
This proves the induction directly by differencing.

Put $\tau=\log_MF$. The powers of $M$ from derivative orders $3,4,5$
are respectively
\[
 (\tfrac12+\tfrac\tau6,\ 1-\tfrac\tau6),\qquad
 (\tfrac57+\tfrac\tau{14},\ \tfrac{29}{28}-\tfrac\tau{14}),\qquad
 (\tfrac56+\tfrac\tau{30},\ \tfrac{25}{24}-\tfrac\tau{30}).
\]
Both entries are at most $5/6+\tau/30$ on, respectively,
$[1,5/2]$, $[5/2,25/8]$, and $[25/8,5]$.
For $\tau\ge5$ the desired bound is trivial. These intervals cover
all $F\ge M$. The argument uses the scale $M$ even on shorter
subintervals, so partial summation gives the same estimate for the
fixed smooth amplitudes in our kernel. \cite[Proposition~\ref*{six-prop:exactphaseproof}]{Weighted}
and \cite[Proposition~\ref*{six-prop:tuplemorseproof}]{Weighted} give the required derivatives
with relative error $O_d(|\eta|)$, for these finitely many orders.
\end{proof}


\section{The thin-energy input and the ten terms}\label{sec:thininput}
We use the curved model of \cite[Section 3]{Weighted}.
The integer-chord count, planar energy, and stationary support are
proved in \cite[Sections 2--3]{Thin}. The weighted thin-energy bound
is \cite[Proposition~\ref*{p-prop:pairprofile}]{Thin}, with proof in
\cite[Appendix~\ref*{p-app:pairproof}]{Thin} of that paper.
These statements hold for $0<k<1/2$ on the stated model ranges.
We use them with the pair proved in Section~\ref{app:elementarypair}.

\begin{proposition}[the ten terms with two finite-derivative pairs]
\label{prop:tentermreturn}
Suppose $(k_p,\ell_p)$ and $(k_d,\ell_d)$ satisfy the finite-derivative
kernel conditions, and set $a_p=1-\ell_p+k_p$, $a_d=1-\ell_d+k_d$,
\[
 \lambda=(T/N)^{k_p}(W/N)^{-a_p},\quad
 \mu=W^{k_d}(W/N)^{-a_d},\quad \theta=k_p+1/2.
\]
In addition to the phase and reduction conditions of \cite{Weighted} assume
\cite[(\ref*{p-eq:thinrange})]{Thin} when $B\ge1$, and
$\lambda B^\theta\le\mu\sqrt U$.
Then the large-value bound for the zeta function consists of the eight terms
$B_1,B_2,B_3,B_4,B_5,B_7,B_8,B_9$ in \cite[(\ref*{six-eq:nine})]{Weighted}, with $k_p$
in the pair terms and the above $\lambda,\mu$, and the two new terms
\begin{equation}\label{eq:newcurvatureterms}
 B_6^{(0)}=\lambda^{1/4}T^{1/16}N^{21/16}V^{-2},\qquad
 B_{10}=\left(\lambda^2T^{1/2+2k_p}N^{21/2}
                         V^{-16-12k_p}\right)^{1/(8+k_p)}.
\end{equation}
\end{proposition}
All terms refer to the same large-value set and the same clustering.
\begin{proof}
Apply \cite[Lemma~\ref*{p-lem:thincompatibility}]{Thin} before the signed $m$ estimate
in each fixed block of \cite[Proposition~\ref*{six-prop:dualityproof}]{Weighted}.
For $D\le B$ use \cite[(\ref*{p-eq:pairprofile})]{Thin}; for $D>B$ the factor $1+B/D$
in \cite[Proposition~\ref*{six-prop:allenergy}]{Weighted} is bounded.
The simultaneous kernel bounds of \cite[Corollary~\ref*{six-cor:two}]{Weighted} therefore give
\[
 \cJ\prec (N/Z)^4\lambda(A_{k_p}(S)+\mathcal C_{k_p}(S;B))
       +(N/Z^2)\mu r^3\sqrt U+rE_2(S;H_0),
 \qquad H_0=C\sqrt{WN^{1/8}}.
\]
The diagonal in \cite[(\ref*{p-eq:pairprofile})]{Thin} is absorbed by the stated
$\lambda B^\theta$ condition. When $B<1$ omit $\mathcal C_{k_p}$.
Both signs of $m$, both Sobolev classes, and all fixed neighbors
are summed as in \cite[Corollary~\ref*{six-cor:blockconsumer}]{Weighted}; their inner
blocks partition the set $S$.

The same thin-energy estimate also covers the small modes.
Their pointwise weight is $D^{-1/2}\le D^{k_p-1/2}$ and their
prefactor divided by $\lambda$ is at most $N^{(a_p-2)/8}$.
This is smaller than one for the pairs used here.
The small-mode diagonal comparison has exponent
$(a_d-2-9k_d)/8<0$. Thus the same energy and absorption apply.
The $F<X$ portion keeps the bound $rE_2(S;H_0)$ in its
local coordinates; the zero mode, nonstationary pair modes and
completion keep their bounds from \cite{Weighted}. The dual coefficients
remain fixed across $m$.

Let $R_1$ be the total number of selected points, $F_c$ the
cluster multiplicity, and $K\prec T/(V^2W)$ the number of classes.
\cite[Proposition~\ref*{six-prop:pairproof}]{Weighted} and the middle-term absorption of
\cite[Lemma~\ref*{six-lem:weighted}]{Weighted}
give simultaneously, for all $1\le H\ll U$,
\[
 \sum_iE_2(S_i;H)\prec R_1\min(H,C_g),\quad
 C_g=1+A_g,\quad A_g=T/(F_c^{3/2}V^6\sqrt{R_1}).
\]
Hence $\sum_iA_{k_p}(S_i)\prec R_1C_g^\theta$ and
\[
 \sum_i\mathcal C_{k_p}(S_i;B)
 \prec R_1\sqrt B\sum_{H\lesssim B}H^{k_p-1}\min(H,C_g)
 \ll R_1\sqrt B\min(B,C_g)^{k_p}.
\]
The last dyadic sum is evaluated below and above $H=C_g$; both
geometric sums are bounded at this transition. This is a simultaneous
bound on the same classes.

The branch of
$r^2Z^2\prec N^2r+NW^{3/4}r^{3/4}\cJ^{1/4}$ given by the term with
$\lambda$, summed by H\"older, gives
\[
 R_1\prec K\lambda^{1/4}N^2W^{3/4}Z^{-3}
                 (C_g^\theta+\sqrt B C_g^{k_p})^{1/4}.
\]
The first part gives $B_2,B_5$. In the second part use
$B\asymp N^{1/2}W^2T^{-3/2}$ and the clustering bounds
$\cN(T,V)\prec ZR_1(TN)^{-1/4}$, $Z\succ F_c(TN)^{1/4}$.
The constant part of $C_g$ gives $F_c^{-2}B_6^{(0)}$.
For its other part, first move $R_1^{-k_p/8}$ to the left:
\[
 R_1^{1+k_p/8}\prec K\lambda^{1/4}N^2W^{3/4}Z^{-3}B^{1/8}
                         (T F_c^{-3/2}V^{-6})^{k_p/4}.
\]
Raise the resulting count to the power $8+k_p$.
Its $Z$ power is $k_p-16<0$, so the lower bound for $Z$ gives
exactly $B_{10}$ in \eqref{eq:newcurvatureterms} and cluster power
$((k_p-16)-3k_p)/(8+k_p)=-2$.
The remaining branches are the unchanged diagonal and small-$F$
calculations in \cite[Section~\ref*{six-sec:return}]{Weighted}; they give the other six listed
terms with nonpositive cluster powers. This proves the proposition.
\end{proof}

\section{The sixth-power upper-amplitude line}
\label{app:upperthree}

\begin{proposition}\label{prop:upperthree}
Using the proved estimate \cite[(\ref*{six-eq:baseatkinson})]{Weighted} and assuming only
\cite[(\ref*{six-eq:basefourth})]{Weighted},
\[
 \cN(T,T^\sigma)\prec T^{E_{\mathrm{up}}(\sigma)},\qquad
 E_{\mathrm{up}}(\sigma)=\frac{242-1340\sigma}{157},\qquad
 \frac{31801}{242699}\le\sigma\le\frac9{68}.
\]
\end{proposition}
This proves the line $(242-1340\sigma)/157$ quoted in
\cite[Section 1]{Thin}, on the stated range, from the displayed inputs.
\begin{proof}
The model estimate is the case $k=3$ of \cite[Theorem 5.1]{dFP}, whose
index $k$ is our $j$ below;
we prove it here using our exact oscillatory integral identities.
We first give the analytic estimate for the straight model
$D_0(w)=\sum a_ne(w\xi_n)$, keeping its cubic frequency term.
Assume $N\le W\le N^{4/3}$ and $|\eta|N^2\ll W$.
Then $U=1+|\eta|W\ll(W/N)^2$ and $\eta^2W\ll1$.
The signed-block proof in \cite[Appendix~\ref*{six-app:dualityproof}]{Weighted} works with
H\"older exponent $2j$, for any fixed $j\ge2$. Its external factor
before the tuple integral is
\[
 N^{3/2+1/(2j)}M^{1/2}r_b^{1-1/(2j)}
\]
in a block with $r_b$ outer points. This follows directly by replacing
four by $2j$ in that proof: the factors are
$N^{3/2}M^{-1/2}(Mr_b)^{1-1/(2j)}(NM)^{1/(2j)}$.
The same two Sobolev classes suffice; their derivative amplitudes
remain bounded. The dual coefficients $b_w$ remain unchanged.

For the straight model, the exact critical phase from
\cite[Proposition~\ref*{six-prop:exactphaseproof}]{Weighted}, with $\kappa=0$, is
\[
 \Theta_m(z,w)=\frac{(z-w)^2}{4Nm}
          +\frac{\eta(z-w)^4}{8(Nm)^3}+\mathcal R_m(z,w).
\]
On the stationary support, all normalized derivatives of $\mathcal R$
are $O_j(\eta^2N|m|)$. Since $N|m|\ll W$ and $\eta^2W\ll1$,
its exponential can be included in a bounded-derivative amplitude.
This step keeps the remainder exactly inside the amplitude.
Write $P=\sum\epsilon_iw_i$, $\Delta_h=\sum\epsilon_iw_i^h$ for
the $2j$ tuple, with $j$ positive and $j$ negative signs.
Set $z=Nmx$. Discard only the phase independent of $x$ after taking
the absolute value of this individual integral. The remaining phase is
\[
 -\frac12\left(P+\frac{\eta\Delta_3}{N^2m^2}\right)x
       +\frac{3\eta\Delta_2}{4Nm}x^2-\frac{\eta P}{2}x^3.
\]
All $w_i$ are in a fixed translated block of length $O(NM)$.
If $|P|\gg U$, its first derivative dominates all the other terms;
repeated integration by parts gives rapid decay. If $|P|\lesssim U$,
the cubic term has bounded derivatives because
$|\eta P|\ll|\eta|+\eta^2W\ll1$, so it too is absorbed exactly
into the amplitude.

For a quadratic $\alpha x+\beta x^2$ on this fixed compact support,
the integral bound is $O((1+|\alpha|+|\beta|)^{-1/2})$.
Indeed, if $|\beta|$ is a fixed fraction of $|\alpha|+|\beta|$,
use the second derivative integral estimate; otherwise the first
derivative is bounded away from zero and one integration by parts
suffices. Thus, with $c_3=\eta\Delta_3/N^2$,
\[
 |I_m|\ll (1+|P+c_3/m^2|)^{-1/2}.
\]
For $|P|\ge1$, either $|c_3|/M^2$ lies outside a fixed comparison range around $|P|$, in
which case the right side is $O(|P|^{-1/2})$ on each of a bounded
number of intervals, or
$|(P+c_3/m^2)'|\asymp|P|/M$.
In the latter case monotonicity and integration between consecutive
integers give
\[
 \sum_{|m|\asymp M}|I_m|\ll 1+M/\sqrt{|P|}.
\]
This also follows by splitting at the possible zero and comparing the
remaining terms with the integral of
$(1+(|P|/M)|x-x_0|)^{-1/2}$.
For $|P|\le1$ the bound $O(M)$ suffices.
Both signs of $m$ have the same estimates. Taking absolute values
here gives the required bound by summing the individual integral estimates.

Put
$J_{2j}=\sum_{1\le D\lesssim U}D^{-1/2}E_{2j}(S;CD)$.
The block tuple sum is bounded by
$(M+\sqrt U)J_{2j}$, including the $O(1)$ term above; tails with
$D>U$ have arbitrarily fast decay. Sum the disjoint inner blocks
and apply H\"older to their outer cardinalities.
As $\sqrt U\ll W/N=:\Mst$,
\[
 M^{j}(M+\sqrt U)\ll \Mst^{j+1}\quad(1\le M\lesssim\Mst).
\]
The full duality bound, including all small modes, is hence
\[
 r^2Z^2\prec N^2r+
 NW^{1/2+1/(2j)}r^{1-1/(2j)}J_{2j}^{1/(2j)}.
\]
This accounts explicitly for the endpoint $1$ in the $m$ sum.

Use $j=3$. \cite[Proposition~\ref*{six-prop:allenergy}]{Weighted} with $B=0$ gives
\[
 J_6\prec (N/Z)^8 A_0(S)+(N/Z^2)r^5\sqrt U,
 \qquad A_0(S)=\sup_D D^{-1/2}E_2(S;D).
\]
Solving the two branches of the preceding inequality gives
\[
 r\prec N^2/Z^2+N^2W^{4/7}Z^{-20/7}A_0(S)^{1/7}
                  +N^{7/2}W^2Z^{-7}U^{1/4}.
\]
Summing the same $K$ classes as in the zeta transfer, and using
$\sum A_0(S_i)\prec R_1C_g^{1/2}$, now gives
\[
 R_1\prec K\left(N^2/Z^2+
 N^{7/3}W^{2/3}Z^{-10/3}C_g^{1/12}
 +N^{7/2}W^2Z^{-7}U^{1/4}\right).
\]

We explain why this straight model applies to the curved sum
in this range. In a height block the difference is
$r_0(w)(n/N)^{3/2}$. If
$B\asymp N^{1/2}W^2T^{-3/2}\ll1$, its exponential has bounded
variation in $n/N$. Abel summation followed by
\cite[Lemma~\ref*{six-lem:completionproof}]{Weighted} therefore gives the straight model,
with the same summed large-value bound and $Z\succ F_c(TN)^{1/4}$.
The summed lower bound selects one common Fourier index, followed
by the selected set and its amplitude level. The clustering proof only requires
$1/10<\sigma<1/6$ and the parameter conditions verified below;
it is therefore valid on the present interval as well.

Substitute $K\prec T/(V^2W)$ and use
$\cN(T,V)\prec ZR_1(TN)^{-1/4}$.
The four terms, before expanding $U$, are
\begin{align*}
 A_1&=T^{1/2}N^{3/2}W^{-1}V^{-2},&
 A_2&=T^{1/6}N^{3/2}W^{-1/3}V^{-2},\\
 A_3&=T^{-3/4}N^{7/4}WV^{-2}(1+NW/T)^{1/4},&
 A_4&=(T^6N^{36}W^{-8}V^{-60})^{1/25}.
\end{align*}
Their cluster powers are respectively
$-1,-7/3,-6,-58/25$, all nonpositive. For example, the $A_g$ branch
has $R_1^{25/24}$ on the left. Raising the resulting count to the power $25$
produces
\[ Z^{-55}(TN)^{-25/4}K^{24}N^{56}W^{16}T^2F_c^{-3}V^{-12}. \]
Insertion of the lower bound for $Z$ proves $A_4$
and its stated multiplicity power.

Put $\nu_*=1-4\sigma$, $\omega_3=(68+184\sigma)/157$, and
$W=N^{q_3}$ with $q_3=\omega_3/\nu_*$. The five exponents
(after expanding $A_3$) at $\nu=\nu_*$ are
\[
 2-8\sigma-\omega_3,\quad 5/3-8\sigma-\omega_3/3,\quad
 1-9\sigma+\omega_3,\quad 1-10\sigma+5\omega_3/4,\quad
 (42-204\sigma-8\omega_3)/25.
\]
The last two equal $E_{\mathrm{up}}$; the first three are smaller throughout the
stated interval. This follows by evaluating their affine differences
at the endpoints: at $\sigma=31801/242699$ the three differences
are $14072/242699$, $139/728097$, $3796/242699$; at $9/68$ they are
$1/17$, $0$, $1/68$.
All five exponents have positive $\nu$ slopes for fixed $q_3$ on the
whole interval, so every $T^{1/3}\le N\prec T^{\nu_*}$ is covered.
Also $9/8<q_3<4/3$ and the exponents of $T$ in
\[
 N/T^{3/5},\quad N^{1/2}W^2/T^{3/2},\quad W/N^{4/3},\quad
 |\eta|N^2/W,\quad |\kappa|W,\quad\eta^2W,\quad
 W\sqrt{N/T}V^{-4}
\]
are strictly negative at both endpoints and throughout all shorter
lengths. The exact values are recorded in the accompanying rational
check; they are affine at $\nu_*$ and have the appropriate monotonicity
in $\nu$. The first pair-count middle term is bounded by
$T^{1/3}<V^4$. Small logarithmic excesses in $N$ fit inside these
strict margins. This proves the asserted line from the two zeta inputs.
\end{proof}

\raggedbottom
\section{All amplitudes with the elementary pair}
\label{app:elementaryassembly}

We complete the proof of Theorem~\ref{thm:elementaryeighth} and
record exactly the cost of avoiding deep exponent pairs.
Use $k_p=k_d=1/30$, $a_p=a_d=1/6$, as proved in
Lemma~\ref{lem:elementaryderivatives}. For $N=T^\nu$, $W=T^\omega$,
write
$\tilde\lambda=\log_T\lambda=k_p+(a_p-k_p)\nu-a_p\omega$ and
$\tilde\mu=\log_T\mu=a_d\nu+(k_d-a_d)\omega$.
The ten exponents in Proposition~\ref{prop:tentermreturn} are
\begin{equation}\label{eq:tencosts}
\begin{aligned}
 e_1&=1/2+3\nu/2-\omega-2\sigma,&
 e_2&=\tilde\lambda/4+1/4+5\nu/4-\omega/4-2\sigma,\\
 e_3&=\tilde\mu/2-1/4+5\nu/4+\omega/2-2\sigma,&
 e_4&=\tilde\mu/2-1/2+3\nu/2+3\omega/4-2\sigma,\\
 e_5&=\frac{4\tilde\lambda+6+4k_p+20\nu-4\omega-(44+24k_p)\sigma}{17+2k_p},&
 e_6&=\tilde\lambda/4+1/16+21\nu/16-2\sigma,\\
 e_7&=1/3+2\nu/3-2\sigma,&
 e_8&=1/2+25\nu/48+\omega/6-10\sigma/3,\\
 e_9&=4(1+\nu-6\sigma)/7,&
 e_{10}&=\frac{2\tilde\lambda+1/2+2k_p+21\nu/2-(16+12k_p)\sigma}{8+k_p}.
\end{aligned}
\end{equation}
Define $\omega(\sigma)$ and $E(\sigma)$ by affine interpolation between
\[
 (\sigma_L,\omega_L,E_L)=(1/8,5/8,481/1024),\qquad
 (\sigma_R,\omega_R,E_R)=\frac1{242699}(31801,153305,102674).
\]
At the right endpoint $e_4=e_5=E_{\mathrm{up}}(\sigma_R)=E_R$.
For every length take $W=N^{q(\sigma)}$, where
$q=\omega(\sigma)/(1-4\sigma)$. Then
$5/4\le q\le30661/23099<4/3$.
Every exponent in \eqref{eq:tencosts} has positive $\nu$ slope
throughout this interval. At $\nu_*=1-4\sigma$ its difference from
$E(\sigma)$ is affine. Here are the endpoint margins, in term order:
\[
\begin{array}{c|cc}
 i&E_L-e_i(\sigma_L)&E_R-e_i(\sigma_R)\\\hline
1&97/1024&24989/242699\\
2&31/15360&1112/3640485\\
3&33/1024&26101/970796\\
4&1/1024&0\\
5&0&0\\
6&31/15360&4261/29123880\\
7&163/3072&25139/728097\\
8&67/3072&77921/11649552\\
9&295/7168&1586/54803\\
10&465/246784&7410/58490459
\end{array}
\]
Thus $\cN(T,T^\sigma)\prec T^{E(\sigma)}$ on this entire interval,
provided the model conditions hold. We verify them explicitly.
The endpoint exponents are
\[
\begin{array}{c|cc}
 \text{ratio}&\sigma_L&\sigma_R\\\hline
 W/N^{4/3}&-1/24&-2065/728097\\
 WB/N^{3/2}&-1/8&-39257/485398\\
 B/N^{1/3}&-1/6&-114568/728097\\
 |\kappa|WB&-1/8&-25483/242699\\
 B^2/W&-5/8&-152687/242699\\
 \lambda B^\theta/(\mu\sqrt U)&-1/15&-6656/117435\\
 |\kappa|W&-1/8&-832/7829\\
 |\eta|N^2/W&-1/8&-49519/242699\\
 |\eta|N^{1/8}&-7/16&-902137/1941592\\
 N^{5/2}T^{-3/2}&-1/4&-75311/242699\\
 \eta^2W&-3/8&-101103/242699\\
 W\sqrt{N/T}V^{-4}&-1/8&-37501/242699
\end{array}
\]
Except for the absorption ratio, these exponents increase with $\nu$
for fixed $q$, or are $(q-4/3)\nu<0$.
For absorption with $B\ge1$, one has $U\asymp T^{\nu+\omega-1}$;
indeed $B\ge1$ and $\nu\le1/2$ imply $\nu+\omega-1>0$.
The exponent of the ratio is bilinear on
$(q,\nu)\in[5/4,30661/23099]\times[1/3,1/2]$.
Its four corner values are
$-2/15,-1/15,-41432/346485,-3190/69297$, all negative.
The small-$M$ comparisons in Proposition~\ref{prop:tentermreturn}
are negative, and $1/3-4\sigma_L<0$ bounds the other middle pair term.
This proves all length and phase conditions with fixed margins.

For the low-amplitude range, the proof of
\cite[Proposition~\ref*{six-prop:curvatureline}]{Weighted} extends verbatim down to
$\sigma_c=27/232$, where
$1-4\sigma_c=(94-500\sigma_c)/67$.
Its four displayed margins remain nonnegative on
$[\sigma_c,5/38]$, its five displayed parameter exponents remain strictly
negative, and $9/8<q\le4/3$ there. Hence for
$\sigma\le\sigma_c$ use $M_4\prec T$, and for
$\sigma_c\le\sigma\le1/8$ use that curvature-only line.
Their eighth-moment contribution is at most $T^{197/134+\varepsilon}$:
$1+4\sigma$ increases up to the first intersection and
$(94-500\sigma)/67+8\sigma$ increases up to $1/8$.
Both classical moments used here follow from the single sparse input
by \cite[Proposition~\ref*{six-prop:momentbaseproof}]{Weighted}.

On $[1/8,\sigma_R]$, $E(\sigma)+8\sigma$ increases and reaches
\[
 E_R+8\sigma_R=\frac{357082}{242699}.
\]
On $[\sigma_R,9/68]$ use Proposition~\ref{prop:upperthree}.
Its eighth-moment line has slope $-84/157<0$.
For $|\zeta|\ge T^{9/68}$ use
$|\zeta|^8\le T^{-4(9/68)}|\zeta|^{12}$ and $M_{12}\prec T^2$,
giving exponent $25/17$. The margins below the exponent
$357082/242699$ for the low range, left middle endpoint, and twelfth-moment tail are
respectively
\[
 37285/32521666,\qquad 389973/248523776,\qquad 2919/4125883.
\]
To pass from separated counts to measure, partition the height axis
into unit cells and select one point above the threshold from each
occupied cell. Three residue classes give separated sets, so their
occupied measure is bounded by three times the separated count.
Use a finite amplitude mesh with all the displayed intersections as
endpoints; the lower-edge count and upper-edge eighth power cost at
most the preassigned $T^\varepsilon$ loss. Below unit amplitude use
the interval length; above $T^{9/68}$ use the preceding integral
inequality directly. Sum dyadic heights, including bounded heights.
This proves Theorem~\ref{thm:elementaryeighth} on the entire range.

\begin{remark}[comparison with the eighth-moment paper]
The main result of \cite{Thin} has exponent $76024973/51678924$. It uses
a Bourgain-derived pair $(k_p,a_p)=(13/414,34/207)$, a Heath-Brown pair
$(k_d,a_d)=(1/56,31/280)$, and the upper-amplitude line
\cite[(\ref*{p-eq:upper})]{Thin} of de Faveri and Pandey.
The ten-term proof also applies to those pairs. Replacing the pairs by
Lemma~\ref{lem:elementaryderivatives} and the line by
Proposition~\ref{prop:upperthree} costs exactly
\[
 \frac{357082}{242699}-\frac{76024973}{51678924}
 =\frac{2428617641}{12542423175876}=0.000193632251\ldots .
\]
This theorem uses fewer cited theorems at the displayed cost.
At the intersection used here, the thin-energy bound obtained from
$E_2(S;H)\ll rH$ gives the result: its
curvature exponent lies below $E_R$ by $494/3640485$.
\cite[Proposition~\ref*{p-prop:pairprofile}]{Thin} also gives the bound using the full pair profile.
\end{remark}

\subsection{A sixth-moment consequence}
\begin{corollary}[sixth moment from curvature energy]
\label{cor:curvaturesixth}
Assume only the sparse-fourth-moment estimate \cite[(\ref*{six-eq:basefourth})]{Weighted}
among the cited inputs of this paper. Then
\[
 M_6(T)\ll_\varepsilon T^{143/116+\varepsilon}.
\]
\end{corollary}
This proof uses the directly proved curvature-energy estimates.
\begin{proof}
Section~\ref{app:elementaryassembly} verifies the extension of
\cite[Proposition~\ref*{six-prop:curvatureline}]{Weighted} to $\sigma_c=27/232$,
including every length and model condition. The already proved
Atkinson estimate is its other zeta input. For $\sigma\le\sigma_c$,
the fourth moment gives $1+2\sigma\le143/116$. On
$[\sigma_c,5/38]$ the sixth-moment line is
\[
 (94-500\sigma)/67+6\sigma=(94-98\sigma)/67.
\]
It decreases from $143/116$. For $\sigma\ge5/38$, the twelfth moment
bounds the contribution by $T^{23/19+\varepsilon}$.
Both classical moments follow from the same cited input
\cite[(\ref*{six-eq:basefourth})]{Weighted}. The same separated-count and dyadic assembly
as in Section~\ref{app:elementaryassembly} completes the proof.
The cost relative to the sixth-moment exponent of \cite{Weighted} is exactly
\[
 \frac{143}{116}-\frac{84599}{68640}=\frac{509}{1990560}.
\]
This alternative uses fewer cited theorems at the displayed
exact exponent cost.
\end{proof}
